% ---
% titre: Problème de proportionnalité- la vente de pâtisserie
% theme: FA
% chapitre: Proportionnalité
% sous_chapitre: Résoudre des problèmes
% niveau: [10e]
% type: EVAL
% tags: [proportionnalite,c-multiplication,  c-division, c-nombres-decimaux, p-question-TS]
% difficulte: 1
% corrige: true
% ---

\question[3]
Lors de notre dernière vente de pâtisseries, nous avons gagné 830\,CHF ce qui représente une somme de 41,50\,CHF par élève. 

\vspace{10pt}
Combien d'argent aurait-il fallu gagner en plus afin que le part de chacun s’élève à 50\,CHF ?

\vspace{10pt}{\footnotesize\raggedleft Espace pour ta démarche et tes calculs\par}
\begin{solutionorgrid}[140pt]

\vspace{5pt}
\def\arraystretch{1.5}
\begin{tabularx}{0.5\textwidth}{l|C|C|}
Somme total [CHF]&830&\\
\hline
Somme par élève [CHF]&41,5&50\\
\end{tabularx}
\def\arraystretch{1.5}

\hfill\textbf{(0.5pt pour la construction du tableau)}

\vspace{15pt}$50 \cdot 830 = 41\,500$ \hfill\textbf{(0.5pt)}

\vspace{5pt}$41\,500 : 41,5 = 1000$\hfill\textbf{(0.5pt)}

\vspace{5pt}$1000 - 830 = 170$ CHF\hfill\textbf{(1pt)}

\end{solutionorgrid}

\begin{tcolorbox}[enhanced jigsaw, interior hidden, colframe=box, parbox=false]%
Réponse: 
\sol{\vspace{20pt}}{Il aurait fallu gagner 170.- en plus. \textbf{(0.5pt)}}
\end{tcolorbox}

